\begingroup\scriptsize\setlength{\tabcolsep}{1.5pt}
\begin{tabular}{@{}lrrrrrrp{2.8cm}@{}}
\hline
Scala & $R$ (m) & $P(Pe_r{=}10)$ (W) & $Pe_r$ a 10 mW & $\omega$ (rad/s) & $\eta_{\rm mecc}$ & $\Delta T$ (K) & Nota \\
\hline
1 nm (Cu) & $5{,}0\times10^{-10}$ & $7{,}90\times10^{3}$ & $1{,}27\times10^{-5}$ & $2{,}48\times10^{4}$ & $7{,}01\times10^{-12}$ & 0{,}0486 & $\eta_{\rm int}\approx$ $2{,}0\times10^{-6}$: il fotone non incontra l'oggetto; $Pe_r\ll10$ \\
1 $\mu$m (vaterite) & $1{,}0\times10^{-6}$ & $1{,}72\times10^{-4}$ & 581{,}18 & 142{,}41 & $4{,}02\times10^{-14}$ & 0{,}0598 & finestra: $f=22{,}7$ Hz \\
10 $\mu$m (vaterite) & $1{,}0\times10^{-5}$ & $1{,}52\times10^{-4}$ & 659{,}88 & 0{,}16 & $4{,}57\times10^{-17}$ & 0{,}0598 & $R>2\,\mu$m: fuori dominio del fascio; estrapolazione \\
\hline
\end{tabular}
\endgroup

\par\smallskip\noindent\footnotesize La riga da 1 m e' stata rimossa: e' un'estrapolazione fuori dominio ($R>2\,\mu$m). Anche la riga da 10 $\mu$m e' oltre il dominio del fascio e va letta come estrapolazione. La riga da 1 nm si riferisce al diametro: $R=0{,}5$ nm.
